\subsection{拓扑空间的子空间}

设 A 是拓扑空间 X 的一个子集。我们已把 X 的拓扑在 A 上诱导的拓扑定义为该拓扑在典范单射 \(A \rightarrow X\) 下的逆像（§ 2，第 3 目，例 1）。一个等价的定义如下：

定义 1. 设 A 是拓扑空间 X 的一个子集。X 的拓扑在 A 上诱导的拓扑是这样的拓扑，其中的开集是 X 的开集与 A 的交。赋以这个拓扑的集合 A 称为 X 的一个子空间。

例。有理直线的拓扑在有理整数集 Z 上诱导的拓扑是离散拓扑，因为 Z 与开区间 \(]n - 1/2, n + 1/2[\) 的交是集合 \(\{n\}\)。

由 § 2 第 3 目命题 5（或直接由定义 1）可以看出，若 B ⊂ A ⊂ X，则 X 的子空间 B 与 X 的子空间 A 的子空间 B 恒同（诱导拓扑的传递性）。若 S 是 X 的拓扑的一个子基（相应地，拓扑基）（§ 2，第 3 目，例 3），则它在 A 上的迹 \(\mathfrak{S}_{\mathrm{A}}\) 是在 A 上诱导的拓扑的一个子基（相应地，拓扑基）。

在所有涉及 A 的元素或子集的问题中，务必仔细区分它们作为 X 的点（相应地，子集）的性质与它们作为子空间 A 的点（相应地，子集）的性质。我们将用"在 A 中"、"关于 A"或"相对于 A"这些说法来指后一类的性质（必要时把它们与"在 X 中"、"关于 X"、"相对于 X"诸说法相对照）。

子空间 A 的开集不必在 X 中是开集；为使 A 中的每个开集都在 X 中是开集，必须且只需 A 在 X 中是开集。该条件是必要的，因为 A 在 A 中是开集；而由 \((\mathrm{O}_{\mathrm{II}})\) 与定义 1，它又是充分的。

A 中的闭集是 X 中的闭集与 A 的交（§ 2，第 3 目，例 1）；与上面一样，我们看出，A 中的每个闭集在 X 中是闭集当且仅当 A 在 X 中是闭集。

点 \(x \in A\) 相对于 A 的诸邻域是 x 相对于 X 的诸邻域与 A 的交。x 相对于 A 的每个邻域都是 x 相对于 X 的邻域，当且仅当 A 是 x 在 X 中的一个邻域。

命题 1. 若 A 与 B 是拓扑空间 X 的两个子集，且 \(B \subset A\)，则 B 在子空间 A 中的闭包是 B 在 X 中的闭包 \(\overline{B}\) 与 A 的交。

若 \(x \in A\)，则 x 在 A 中的每个邻域都有 \(V \cap A\) 的形式，其中 V 是 x 在 X 中的一个邻域。由于 \(V \cap B = (V \cap A) \cap B\)，由此得出，x 属于 B 相对于 A 的闭包当且仅当 x 属于 B 相对于 X 的闭包。

推论. A 的子集 B 在 A 中稠密，当且仅当在 X 中 \(\overline{B} = \overline{A}\)（亦即当且仅当 \(A \subset \overline{B}\)）。

由此推出，若 A、B、C 是 X 的三个子集，满足 \(A \supset B \supset C\)，且 B 在 A 中稠密，C 在 B 中稠密，则 C 在 A 中稠密（稠密性的传递性），因为在 X 中有 \(\overline{A} = \overline{B} = \overline{C}\)。

命题 2. 设 A 是拓扑空间 X 的一个稠密子集；则对每个 \(x \in \mathbf{A}\) 以及 \(x\) 相对于 A 的每个邻域 V，V 在 X 中的闭包 \(\overline{\mathbf{V}}\) 是 \(x\) 在 X 中的一个邻域。

因为 V 包含 U ∩ A，其中 U 是 X 的包含 x 的开子集，故 \(\overline{V}\) 包含 U ∩ \(\overline{A} = U\)（§ 1，第 6 目，命题 5）。

命题 3. 设 \((\mathbf{A}_i)_{i\in I}\) 是拓扑空间 \(\mathbf{X}\) 的子集的一个族，使得下列性质之一成立：

\begin{enumerate}
\def\labelenumi{\alph{enumi})}
\item
  诸 \(A_{i}\) 的内部覆盖 X。
\item
  \((\mathbf{A}_{\iota})_{\iota \in \mathbf{I}}\) 是 \(\mathbf{X}\) 的一个局部有限闭覆盖（\(\S 1\)，第 5 目）。
\end{enumerate}

在这些条件下，X 的子集 B 在 X 中是开（相应地，闭）集，当且仅当诸集合 \(B \cap A_{t}\) 中的每一个在 \(A_{t}\) 中都是开（相应地，闭）集。

显然，若 B 在 X 中是开（相应地，闭）集，则 \(B \cap A_{t}\) 在 \(A_{t}\) 中是开（相应地，闭）集。反之，先设条件 a) 满足；由于 \((\mathbb{C}\mathbf{B}) \cap \mathbf{A}_{t} = \mathbf{A}_{t} - (\mathbf{B} \cap \mathbf{A}_{t})\)，由对偶性，只需考虑每个 \(B \cap A_{t}\) 相对于 \(A_{t}\) 为开集的情形。在这一情形，\(B \cap \mathring{A}_{t}\) 在 \(\mathring{A}_{t}\) 中是开集（对每个 \(t \in I\)），从而在 X 中是开集；又由假设有 \(B = \bigcup_{i} (B \cap \mathring{A}_{t})\)，故 B 在 X 中是开集。

现在设 b) 满足；再由对偶性，只需考虑每个 \(B \cap A_{t}\) 在 \(A_{t}\) 中是闭集、从而在 X 中是闭集的情形。由于族 \((B \cap A_{t})\) 局部有限且 \(\mathbf{B} = \bigcup_{t} (B \cap A_{t})\)，由 §1 第 5 目命题 4 可知，B 在 X 中是闭集。

注. 设 \((\mathbf{U}_i)_{i\in I}\) 是拓扑空间 \(\mathbf{X}\) 的一个开覆盖，且对每个 \(i\in I\)，设 \(\mathfrak{B}_i\) 是子空间 \(\mathbf{U}_i\)（\(\mathbf{X}\) 的子空间）的拓扑的一个基；则显然 \(\mathfrak{B} = \bigcup_{i\in I}\mathfrak{B}_i\) 是 \(\mathbf{X}\) 的拓扑的一个基。

